\chapter{绪论}\label{chap:guide}

\section{研究背景}

八皇后问题是典型的约束满足问题（CSP）：变量数量不大、约束关系清晰、又能完整覆盖“建模-搜索-剪枝-验证”这条实践链路。它既适合用命令式回溯实现，也能转写为逻辑关系求解，因此非常适合用于比较不同编程范式在表达力和执行效率上的取舍\cite{russell2021ai,byrd2006minikanren,mackworth1977consistency}。

\section{任务要求拆解}

结合老师的实验要求pdf，本次实验报告可拆分为6个分实验任务：

\begin{enumerate}
  \item 先用非逻辑方法生成候选，再用逻辑关系做合法性筛选；
  \item 给出不依赖命令式全排列的纯逻辑求解方案；
  \item 补充 Kanren 中缺失的不等约束关系并用于八皇后建模；
  \item 使用 DFS 和 BFS 求解，并与逻辑方案做对照；
  \item 记录调试过程，说明出现的问题与修复策略；
  \item 在课程要求基础上给出一条可落地的性能改进路径（本报告采用位运算回溯）。
\end{enumerate}